% !TEX program = pdflatex
\documentclass[aspectratio=169,11pt]{beamer}

% Červený vzhľad
\usetheme{Madrid}
\usecolortheme{beaver}
\setbeamercolor{alerted text}{fg=black!70!red}
\setbeamercolor{frametitle}{bg=red!70!black, fg=white}
\setbeamercolor{title}{bg=red!70!black,fg=white}
\setbeamercolor{block title}{bg=red!70!black,fg=white}
\setbeamercolor{block body}{bg=red!5!white}
\setbeamercolor{itemize item}{fg=red!70!black}
\setbeamercolor{itemize subitem}{fg=red!70!black}
\setbeamercolor{enumerate item}{fg=red!70!black}

\usepackage[utf8]{inputenc}
\usepackage[T1]{fontenc}
\usepackage[slovak]{babel}
\usepackage{tikz}
\usetikzlibrary{positioning,arrows.meta,shapes.geometric,calc,fit}

% TikZ styly pro diagramy
\tikzset{
  component/.style={rectangle,rounded corners=3pt,draw=red!70!black,thick,
                    fill=red!10,minimum width=2.6cm,minimum height=1.0cm,
                    align=center,font=\small},
  store/.style={cylinder,shape border rotate=90,draw=red!70!black,thick,
                fill=red!8,aspect=0.35,minimum width=2.0cm,minimum height=1.6cm,
                align=center,font=\small},
  queue/.style={rectangle,draw=red!70!black,thick,fill=red!15,
                minimum width=2.8cm,minimum height=0.8cm,
                align=center,font=\small\ttfamily},
  arr/.style={-{Stealth[length=2mm]},thick,red!60!black},
  arrdash/.style={-{Stealth[length=2mm]},thick,dashed,red!60!black},
  lab/.style={font=\scriptsize,fill=white,inner sep=1pt},
}

\title[Big Data Pipeline Monitor]{Big Data Pipeline Monitor}
\subtitle{Architektúra distribuovanej aplikácie}
\author{Jozef Waldhauser}
\institute{MSWA}
\date{2026 léto}

\begin{document}

\begin{frame}
  \titlepage
\end{frame}

\begin{frame}{Obsah}
	\tableofcontents
\end{frame}

% ---------------------------------------------------------------
\section{Prehľad}
\begin{frame}{3 komponenty}
  \begin{itemize}
    \item Aplikácia eviduje \textbf{datasety}, \textbf{pipeline}, ich \textbf{behy} a \textbf{alerty}.
    \item Spustenie pipeline je \textbf{simulovaný výpočet}\,-\,reálne sa robí agregácia nad MongoDB.
    \item Architektonicky som riešenie rozdelil na \textbf{tri samostatné služby}, ktoré komunikujú cez REST a frontu správ.
  \end{itemize}
  \vfill
  \begin{block}{Tri komponenty}
    \texttt{qr-frontend} \,$\;\to\;$\, \texttt{qr-backend} \,$\;\to\;$\, \texttt{qr-computation-module}
  \end{block}
\end{frame}

% ---------------------------------------------------------------
\section{Architektúra}
\begin{frame}[fragile]{Architektúra}
  \centering
  \begin{tikzpicture}[node distance=1.0cm and 1.4cm]
	
	\node[component] (fe) {qr-frontend\\\scriptsize Next.js / React};
	
	\node[component, right=3cm of fe] (be) {qr-backend\\\scriptsize Express + Mongoose};
	
	\node[component, right=3cm of be] (wk) {qr-computation\\\scriptsize Node worker};
	
	\node[store, below=1.5cm of be] (db) {MongoDB};
	
	\node[queue, above=1.0cm of be] (mq) {computations queue};
	
	\draw[arr]   (fe) -- node[lab,above]{HTTP / JSON} (be);
	
	\draw[arr]   (be) -- node[lab,left]{publish} (mq);
	
	\draw[arr]   (mq) -- node[lab,right]{consume} (wk);
	
	\draw[arr]   (wk.south) to[bend left=20] node[lab,below]{PATCH /runs/(id)} (be.south east);
	
	\draw[arr]   (be) -- node[lab,right]{Mongoose ODM} (db);
	
	\draw[arrdash] (wk.south) to[bend left=45] node[lab,below]{aggregate auditlogs} (db.east);
	
\end{tikzpicture}
  \begin{itemize}\small
	\item \textbf{Synchrónna cesta:} prehliadač $\to$ API (CRUD a stav).
	\item \textbf{Asynchrónna cesta:} API $\to$ fronta $\to$ worker $\to$ späť do API.
	\item Každá služba má \textbf{vlastný proces a vlastný port}\,---\,nasaditeľné a škálovateľné samostatne.
\end{itemize}

\end{frame}

% ---------------------------------------------------------------
\section{Backend}
\begin{frame}{Backend: \texttt{qr-backend}}
  \begin{itemize}
    \item Express 5, modulárne routy v \texttt{app/}\,:\, \texttt{datasets}, \texttt{pipelines}, \texttt{runs}, \texttt{alert-rules}, \texttt{alerts}.
    \item Perzistencia cez \textbf{Mongoose} (MongoDB)\,-\,jeden model na entitu (\texttt{Dataset}, \texttt{Pipeline}, \texttt{JobRun}, \texttt{AlertRule}, \texttt{Alert}).
    \item \textbf{Validácia} vstupu \textbf{Joi} schémami v každom module (\texttt{validation.js}).
    \item \textbf{Error handling}: per-route \texttt{try/catch}, centrálny 404 a 500 handler v \texttt{server.js}.
    \item Publikácia úloh do RabbitMQ cez \texttt{app/rabbitmq.js} (front \texttt{computations}, \texttt{durable}).
  \end{itemize}
  \begin{block}{Oddelenie zodpovedností}
    Router robí validáciu a HTTP, model drží schému, business pravidlá sú v handleri (napr. kontrola \texttt{active}, stavový prechod).
  \end{block}
\end{frame}

% ---------------------------------------------------------------
\section{Doménový model}
\begin{frame}{Doménový model (Mongoose)}
  \small
  \begin{itemize}
    \item \textbf{Dataset}\,-\,metadáta (\texttt{name}, \texttt{owner}, \texttt{type: aggregation|file}, \texttt{aggregation.timeFrom/To}).
    \item \textbf{Pipeline}\,-\,odkaz na dataset (\texttt{datasetOid}), \texttt{schedule}, \texttt{active}, \texttt{pipelineVersion}, \texttt{lastStatus}, \texttt{lastRunTime}.
    \item \textbf{JobRun}\,-\,konkrétny beh: \texttt{pipelineOid}, \texttt{pipelineVersion}, \texttt{status}, časy a \texttt{processedRecords}.
    \item \textbf{AlertRule}\,-\,definícia (\texttt{pipelineOid}, \texttt{reportWhenState}).
    \item \textbf{Alert}\,-\,vyvolaný záznam (väzba na pipeline + run, \texttt{acknowledgedAt}).
  \end{itemize}
  \begin{block}{Pravidlo integrity}
    Pipeline nemôže vzniknúť nad neexistujúcim datasetom\,-\,kontroluje sa pred uložením v handleri \texttt{POST /pipelines}.
  \end{block}
\end{frame}

% ---------------------------------------------------------------
\section{Computation module}
\begin{frame}{Computation module: \texttt{qr-computation-module}}
  \begin{itemize}
    \item Samostatný Node proces s vlastným Express serverom a RabbitMQ consumerom (\texttt{app/rabbit-consumer}).
    \item Pre každú správu \texttt{calcStats}:
      \begin{enumerate}\small
        \item \texttt{PATCH /runs/:id} $\to$ \texttt{running} (s \texttt{startTime}),
        \item agregácia nad MongoDB kolekciou \texttt{auditlogs} v zadanom časovom intervale,
        \item uloženie výsledku do \texttt{computationstats},
        \item \texttt{PATCH /runs/:id} $\to$ \texttt{successful} / \texttt{error}.
      \end{enumerate}
    \item \textbf{Idempotencia}\,-\,upsert \texttt{ComputationJob} podľa \texttt{runId} + \texttt{receiveCount}.
    \item \textbf{Odolnosť}\,-\,reconnect s exponenciálnym backoffom, \texttt{prefetch(1)}, \texttt{ack}/\texttt{nack}.
  \end{itemize}
\end{frame}

% ---------------------------------------------------------------
\section{Frontend}
\begin{frame}{Frontend: \texttt{qr-frontend}}
  \begin{itemize}
    \item \textbf{Next.js} App Router\,-\,React Server Components načítavajú dáta priamo z API počas renderu (\texttt{src/lib/api.ts}).
    \item Stránky:\,\texttt{/}\,(dashboard), \texttt{/datasets}, \texttt{/pipelines}, \texttt{/pipelines/[id]}, \texttt{/runs}, \texttt{/runs/[id]}, \texttt{/alerts}, \texttt{/alert-rules}.
    \item Klientske komponenty pre akcie a živý stav:\;\texttt{RunPipelineButton}, \texttt{RunStatusLive}, \texttt{PipelineActiveToggle}, \texttt{AcknowledgeAlertButton}.
    \item \textbf{Stavy UI}: \textit{loading} (RSC streaming), \textit{empty} (prázdne zoznamy), \textit{error} (chybové hlášky z \texttt{ApiRequestError}).
    \item Cez \texttt{next.config} rewrite na \texttt{/blocky-api}\,-\,same-origin volania z prehliadača (žiadne CORS).
  \end{itemize}
\end{frame}

% ---------------------------------------------------------------
\section{Pipeline}
\begin{frame}[fragile]{Komunikácia: spustenie pipeline}
  \centering
\begin{tikzpicture}[node distance=0.6cm and 1.7cm]
	\node[component] (ui) {Frontend};
	\node[component, right=of ui] (api) {Backend API};
	\node[queue,     right=of api] (q)   {computations};
	\node[component, right=of q]   (w)   {Worker};
	\node[store, below=1.4cm of api] (m) {MongoDB};
	
	% POST: arc nad oboma uzlami, label neprekrýva boxy
	\draw[arr] (ui.north) to[out=75,in=105]
	node[lab,above]{POST /pipelines/(id)/run} (api.north);
	
	% priame šípky
	\draw[arr] (api) -- node[lab,left]{insert JobRun} (m);
	\draw[arr] (api) -- node[lab,above]{enqueue} (q);
	\draw[arr] (q)   -- node[lab,above]{consume}  (w);
	
	% PATCH: L-cesta NAD riadkom uzlov (vyhne sa q)
	\draw[arr] (w.north) -- ++(0,0.6cm)
	-- node[lab,above]{PATCH /runs/(id)} ($(api.north)+(0,0.6cm)$)
	-- (api.north);
	
	% aggregate: L-cesta POD riadkom uzlov (vyhne sa q)
	\draw[arrdash] (w.south) -- (w.south |- m.east)
	-- node[lab,below]{aggregate} (m.east);
\end{tikzpicture}
  \vfill
  \begin{itemize}\small
    \item API vráti odpoveď \textbf{ihneď} po zaradení do fronty\,-\,žiadne dlhé HTTP požiadavky.
    \item Pri \textbf{nedostupnom RabbitMQ} backend run uloží ako \texttt{error} a vráti \texttt{503}.
  \end{itemize}
\end{frame}

% ---------------------------------------------------------------
\section{Stav pipeline}
\begin{frame}[fragile]{Stavový stroj behu a biznis pravidlá}
  \centering
  \begin{tikzpicture}[node distance=1.3cm,every node/.style={font=\small}]
    \node[component,fill=red!5] (p) {pending};
    \node[component,right=of p,fill=red!15] (r) {running};
    \node[component,above right=0.3cm and 1.3cm of r,fill=green!15,draw=green!50!black] (s) {successful};
    \node[component,below right=0.3cm and 1.3cm of r,fill=red!25,draw=red!70!black] (e) {error};
    \draw[arr] (p) -- (r);
    \draw[arr] (r) -- (s);
    \draw[arr] (r) -- (e);
  \end{tikzpicture}
  \vfill
  \begin{itemize}
    \item Povolené prechody sú validované v \texttt{PATCH /runs/:id}\,-\,iné kombinácie vrátia \texttt{409}.
    \item \textbf{Pipeline} sa dá spustiť len ak \texttt{active = true} (inak \texttt{409}).
    \item \textbf{Alert} vznikne podľa \texttt{AlertRule.reportWhenState} pri prechode behu do daného stavu.
  \end{itemize}
\end{frame}

% ---------------------------------------------------------------
\section{Designové rozhodnutia}
\begin{frame}{Návrhové rozhodnutia}
  \begin{itemize}
    \item \textbf{Tri služby namiesto monolitu}\,-\,API zostáva responzívne, výpočet je odolný voči pomalým úlohám, dá sa nezávisle škálovať.
    \item \textbf{RabbitMQ ako buffer}\,-\,durable fronta, \texttt{prefetch(1)}, manuálne ack/nack\,$\Rightarrow$\,žiadna strata úloh pri reštarte.
    \item \textbf{MongoDB + Mongoose}\,-\,doménové entity sú prirodzene dokumentové, schémy v Mongoose dávajú validáciu a typovú istotu.
    \item \textbf{Joi validácia + jednotné HTTP kódy}\,-\,400 / 404 / 409 / 503 podľa významu.
    \item \textbf{Stavový stroj na serveri}\,-\,nedá sa „preskočiť“ priamo z \texttt{pending} do \texttt{successful}.
    \item \textbf{Next.js RSC} na FE\,-\,výpisy sa renderujú serverom z API, klientské komponenty len pre akcie a polling stavu.
  \end{itemize}
\end{frame}

% ---------------------------------------------------------------
\begin{frame}
  \centering
  \vfill
  {\Large \textbf{Ďakujem za pozornosť.}}\\[1em]
  {\large Otázky?}
  \vfill
\end{frame}

\end{document}
